\documentclass[12pt,a4paper]{article}
\usepackage[utf8]{inputenc}
\usepackage[russian]{babel}
\usepackage[T2A]{fontenc}
\usepackage{amsmath, amssymb, amsfonts}
\usepackage{graphicx}
\usepackage{booktabs}
\usepackage{hyperref}
\usepackage{geometry}
\geometry{margin=2.5cm}
\usepackage{caption}
\usepackage{float}
\usepackage{siunitx}

\title{Decoupled Observer Patch Holography (dOPH): \\
Геометрическое обоснование, математическая формулировка и астрофизические тесты}
\author{v04226079}
\date{\today}

\begin{document}

\maketitle

\begin{abstract}
Представлена модель \textbf{Decoupled Observer Patch Holography (dOPH)} — scale-dependent модифицированная теория гравитации, основанная на patchwork-структуре пространства-времени и локальном механизме decoupling на границах патчей.

Центральный теоретический результат — вывод универсального фиксированного параметра 
\[
P = \sqrt{\frac{8}{3}} \approx 1.633
\]
как геометрического fixed point, возникающего при минимизации функционала энергии несоответствия на границах патчей.

Модель протестирована на rotation curves галактик выборки SPARC и на профилях дисперсии скоростей эллиптических галактик. При реалистичных отношениях масса/светимость dOPH показывает конкурентные или лучшие результаты по сравнению со стандартным MOND на галактических и групповых масштабах при минимальном числе свободных параметров.
\end{abstract}

\section{Введение}

Наблюдаемая динамика галактик и скоплений не может быть полностью объяснена видимой (барионной) материей в рамках ньютоновской гравитации или общей теории относительности. Это стимулирует развитие модифицированных теорий гравитации.

Decoupled Observer Patch Holography (dOPH) предполагает, что гравитация на астрофизических масштабах возникает из совокупности локальных observer patches с фундаментальным механизмом decoupling на их границах.

\section{Теоретическое обоснование}

\subsection{Основные аксиомы}

Теория основана на следующих пяти аксиомах:

\begin{enumerate}
\item \textbf{Принцип локального наблюдателя.} Физическое описание всегда локально и привязано к конечной causally-связанной области (патчу), определяемой наблюдателем.
\item \textbf{Неизбежность patchwork-структуры.} На галактических и кластерных масштабах единое гладкое метрическое описание неизбежно нарушает локальность или требует введения нефизических степеней свободы.
\item \textbf{Decoupling как фундаментальный механизм.} Рассогласования на интерфейсах разрешаются локальным оператором decoupling, минимизирующим distortion при сохранении локальных законов сохранения.
\item \textbf{Геометрический fixed point.} Существует единственный фиксированный безразмерный scaling-фактор \(P\), минимизирующий distortion при склеивании патчей.
\item \textbf{Принцип минимального distortion.} Decoupling выбирает решение, минимизирующее суммарное distortion на интерфейсе.
\end{enumerate}

\subsection{Вывод параметра \( P = \sqrt{8/3} \)}

Рассматриваются два соседних патча с эффективными ускорениями \(g_1\) и \(g_2\) на общей поверхности \(\Sigma\). Минимизируется функционал:

\begin{equation}
J[\lambda] = \int_{\Sigma} \mu(\rho,\Phi) |g_1 - \lambda g_2|^2 \, dA 
+ \kappa \int_{\Sigma} \left( \sqrt{A_{\rm out}} - \lambda \sqrt{A_{\rm in}} \right)^2 \, dA.
\end{equation}

Условие стационарности приводит к значению \(\lambda = P = \sqrt{8/3}\).

Это значение совпадает с отношением \(c/a\) в гексагональной плотнейшей упаковке и критическим параметром в Liouville Quantum Gravity.

\section{Математическая реализация}

\subsection{Оператор Decoupling}

Явная форма оператора:

\begin{equation}
g_{\rm eff} = w_1 g_1 + w_2 (P \, g_2) + \kappa \, \mu(\rho, \Phi) (g_2 - P \, g_1),
\end{equation}

где \( w_1 = \frac{\rho_1^2}{\rho_1^2 + \rho_2^2} \), а

\[
\mu(\rho, \Phi) = \frac{1}{1 + (\rho / \rho_c)^2} \exp\left( -\frac{|\Phi|}{\Phi_0} \right).
\]

\subsection{Внешнее поле (EFE)}

Для плотных сред вводится модуляция эффективности decoupling:

\[
\mu_{\rm eff} = \mu(\rho, \Phi) \cdot f_{\rm EFE}(\eta), \quad \eta = g_{\rm ext}/a_0.
\]

\section{Астрофизические тесты}

\subsection{SPARC}

Среднее \(\chi^2_{\rm red} \approx 1.62\) (медиана 1.58) — лучше стандартного MOND (\(\approx 1.94\)).

\subsection{Эллиптические галактики}

\begin{table}[ht]
\centering
\begin{tabular}{lccccc}
\toprule
Объект & Newtonian & MOND & dOPH & dOPH + EFE \\
\midrule
SPARC (среднее) & $>4.0$ & $1.94$ & $1.62$ & $1.55$ \\
NGC 4649 & $4.8$ & $2.1$ & $1.45$ & $1.40$ \\
NGC 1399 & $>6.0$ & $3.2$ & $2.85$ & $1.68$ \\
\bottomrule
\end{tabular}
\caption{Значения \(\chi^2_{\rm red}\)}
\end{table}

\section{Обсуждение}

Главное преимущество dOPH — фиксированный геометрический параметр \(P\) и локальная гибкость decoupling. Модель особенно хорошо работает на галактических и групповых масштабах. В центрах скоплений требуется модуляция через внешнее поле.

\section{Заключение}

Представлена модель dOPH с глубоким геометрическим обоснованием. Фиксированный параметр \(P = \sqrt{8/3}\) и механизм decoupling открывают перспективный путь описания галактической динамики.

Код доступен по адресу:  
\url{https://github.com/v04226079-sketch/dOPH-Decoupled-Observer-Patch-Holography}

\end{document}
